\documentclass[12pt,letterpaper]{article}

\usepackage[utf8]{inputenc}
\usepackage[T1]{fontenc}
\usepackage[spanish,es-tabla]{babel}
\usepackage{lmodern}
\usepackage{microtype}
\usepackage[margin=2.5cm]{geometry}
\usepackage{setspace}
\usepackage{graphicx}
\usepackage{float}
\usepackage{tabularx}
\usepackage{array}
\usepackage{xcolor}
\usepackage{listings}
\usepackage{enumitem}
\usepackage[hidelinks]{hyperref}

\setstretch{1.25}
\setlength{\parindent}{0pt}
\setlength{\parskip}{0.55em}
\sloppy
\renewcommand{\arraystretch}{1.25}

\definecolor{codegray}{RGB}{245,245,245}
\definecolor{codeframe}{RGB}{180,180,180}
\definecolor{keywordcolor}{RGB}{120,40,130}
\definecolor{commentcolor}{RGB}{45,110,60}
\definecolor{stringcolor}{RGB}{170,70,30}

\lstdefinestyle{codigo}{
    backgroundcolor=\color{codegray},
    frame=single,
    rulecolor=\color{codeframe},
    basicstyle=\ttfamily\small,
    keywordstyle=\color{keywordcolor}\bfseries,
    commentstyle=\color{commentcolor},
    stringstyle=\color{stringcolor},
    numbers=left,
    numberstyle=\tiny,
    numbersep=8pt,
    breaklines=true,
    breakatwhitespace=false,
    showstringspaces=false,
    tabsize=4,
    columns=fullflexible,
    keepspaces=true,
    captionpos=b
}

\lstset{style=codigo}
\newcolumntype{Y}{>{\centering\arraybackslash}X}

\newcommand{\evidencia}[3]{
\begin{figure}[H]
\centering
\IfFileExists{#1}{
    \includegraphics[width=0.95\textwidth]{#1}
}{
    \fbox{
        \begin{minipage}[c][6cm][c]{0.90\textwidth}
        \centering
        \textbf{#2}
        \end{minipage}
    }
}
\caption{#3}
\end{figure}
}



\begin{document}

% ============================================================
%  SECCIÓN 4.4 — Pruebas data-driven (Diego Calva)
% ============================================================
\section{Pruebas data-driven y conexión de cobertura con \texttt{pytest-cov}}

\subsection{Objetivo de la implementación}

Para reducir la duplicación de código y ejecutar un mismo flujo con distintos datos de entrada, se diseñó una prueba \textit{data-driven} para el módulo de autenticación de Cal.com. La implementación utiliza el decorador \texttt{@pytest.mark.parametrize} y obtiene los datos desde un archivo CSV externo.

La separación entre la lógica de prueba y los datos permite ampliar los escenarios sin crear una función diferente para cada combinación. Este enfoque también mejora la mantenibilidad de la suite y facilita la identificación de fallas específicas.

La cobertura se integró mediante \texttt{pytest-cov}, con salida en terminal, reporte HTML y archivo XML compatible con SonarQube.

\subsection{Diseño de los casos de prueba}

La prueba valida el inicio de sesión de Cal.com con datos válidos, inválidos y de frontera. Los escenarios definidos se muestran en la Tabla~\ref{tab:data-driven-login}.

\begin{table}[H]
\centering
\caption{Casos data-driven para el inicio de sesión de Cal.com}
\label{tab:data-driven-login}
\begin{tabularx}{\textwidth}{|>{\centering\arraybackslash}p{1.4cm}|X|X|}
\hline
\textbf{ID} & \textbf{Datos de entrada} & \textbf{Resultado esperado} \\
\hline
DD-01 & Correo y contraseña válidos & El usuario inicia sesión y abandona la ruta \texttt{/auth/login}. \\
\hline
DD-02 & Correo sin formato válido & El navegador muestra una validación del campo de correo. \\
\hline
DD-03 & Correo vacío & El formulario marca el campo de correo como obligatorio. \\
\hline
DD-04 & Contraseña vacía & El formulario marca el campo de contraseña como obligatorio. \\
\hline
DD-05 & Contraseña incorrecta & La aplicación muestra un error de autenticación y no inicia sesión. \\
\hline
\end{tabularx}
\end{table}

Las credenciales válidas se obtienen mediante las variables de entorno \texttt{CALCOM\_TEST\_EMAIL} y \texttt{CALCOM\_TEST\_PASSWORD}. Con ello, el repositorio puede conservar los casos de prueba sin exponer información sensible.

\subsection{Archivo de datos CSV}

Los datos de entrada se almacenaron en el archivo \texttt{tests/data/login\_cases.csv}. Cada fila representa una ejecución independiente de la misma prueba.

\begin{lstlisting}[caption={Archivo CSV con los casos de prueba data-driven}]
id,email,password,expected,description
DD-01,ENV:CALCOM_TEST_EMAIL,ENV:CALCOM_TEST_PASSWORD,success,Credenciales validas
DD-02,correo-sin-arroba,Password123!,validation,Formato de correo invalido
DD-03,,Password123!,validation,Correo vacio
DD-04,usuario@example.com,,validation,Contrasena vacia
DD-05,usuario@example.com,ContrasenaIncorrecta123!,auth_error,Contrasena incorrecta
\end{lstlisting}

El prefijo \texttt{ENV:} identifica los valores que se recuperan desde variables de entorno. Esta estructura mantiene separados los datos sensibles y los datos versionados.

\subsection{Implementación de la prueba parametrizada}

La función \texttt{load\_login\_cases()} abre el archivo CSV, transforma cada fila en un diccionario y sustituye los valores con prefijo \texttt{ENV:} por el contenido de las variables de entorno correspondientes.

\begin{lstlisting}[language=Python,caption={Prueba data-driven con \texttt{@pytest.mark.parametrize}}]
from __future__ import annotations

import csv
import os
from pathlib import Path

import pytest

from pages.login_page import LoginPage

DATA_FILE = Path(__file__).parent / "data" / "login_cases.csv"


def _resolve_value(value: str) -> str:
    if value.startswith("ENV:"):
        return os.getenv(value.removeprefix("ENV:"), "")
    return value


def load_login_cases() -> list[dict[str, str]]:
    with DATA_FILE.open(encoding="utf-8", newline="") as file:
        cases = list(csv.DictReader(file))

    for case in cases:
        case["email"] = _resolve_value(case["email"])
        case["password"] = _resolve_value(case["password"])
    return cases


@pytest.mark.e2e
@pytest.mark.parametrize(
    "case",
    load_login_cases(),
    ids=lambda case: case["id"],
)
def test_login_data_driven(
    driver,
    base_url: str,
    case: dict[str, str],
) -> None:
    if case["id"] == "DD-01" and (
        not case["email"] or not case["password"]
    ):
        pytest.skip(
            "Credenciales de prueba no disponibles "
            "en las variables de entorno."
        )

    login_page = LoginPage(driver, base_url)
    login_page.open()
    login_page.submit_credentials(
        case["email"],
        case["password"],
    )

    result = login_page.wait_for_result()

    assert result.status == case["expected"], (
        f"Caso {case['id']} ({case['description']}): "
        f"se esperaba '{case['expected']}', "
        f"pero se obtuvo '{result.status}'. "
        f"Detalle: {result.message}"
    )
\end{lstlisting}

Aunque existe una sola función, pytest registra cada fila del archivo CSV como una ejecución separada. Los casos \texttt{DD-01}, \texttt{DD-02}, \texttt{DD-03}, \texttt{DD-04} y \texttt{DD-05} aparecen individualmente en la terminal y en el reporte HTML.

El argumento \texttt{ids=lambda case: case["id"]} asigna a cada ejecución el identificador definido en el CSV. Esto mejora la lectura de los resultados y permite localizar con rapidez el conjunto de datos relacionado con una falla.

\evidencia
{imagenes/salida_pytest.png}
{Evidencia de ejecución de los casos data-driven en pytest}
{Salida de pytest con la ejecución individual de los casos DD-01 a DD-05.}

\subsection{Integración con Page Object Model}

La prueba se apoya en el Page Object \texttt{LoginPage}. Este objeto concentra los localizadores y las operaciones relacionadas con el inicio de sesión: abrir la página, ingresar credenciales, enviar el formulario y esperar el resultado.

\begin{lstlisting}[language=Python,caption={Fragmento principal del Page Object \texttt{LoginPage}}]
from dataclasses import dataclass

from selenium.common.exceptions import TimeoutException
from selenium.webdriver.common.by import By
from selenium.webdriver.support import expected_conditions as EC
from selenium.webdriver.support.ui import WebDriverWait


@dataclass(frozen=True)
class LoginResult:
    status: str
    message: str = ""


class LoginPage:
    EMAIL = (
        By.CSS_SELECTOR,
        'input[name="email"], input[type="email"]',
    )
    PASSWORD = (
        By.CSS_SELECTOR,
        'input[name="password"], input[type="password"]',
    )
    SUBMIT = (By.CSS_SELECTOR, 'button[type="submit"]')
    ERROR = (
        By.CSS_SELECTOR,
        '[role="alert"], [data-testid="toast-error"], '
        '[data-testid*="error"]',
    )

    def __init__(self, driver, base_url: str, timeout: int = 12):
        self.driver = driver
        self.base_url = base_url.rstrip("/")
        self.wait = WebDriverWait(driver, timeout)

    def open(self) -> None:
        self.driver.get(f"{self.base_url}/auth/login")
        self.wait.until(
            EC.visibility_of_element_located(self.EMAIL)
        )

    def submit_credentials(self, email: str, password: str) -> None:
        email_input = self.wait.until(
            EC.visibility_of_element_located(self.EMAIL)
        )
        email_input.clear()
        email_input.send_keys(email)

        password_fields = self.driver.find_elements(*self.PASSWORD)
        if not password_fields:
            self.wait.until(
                EC.element_to_be_clickable(self.SUBMIT)
            ).click()
            password_input = self.wait.until(
                EC.visibility_of_element_located(self.PASSWORD)
            )
        else:
            password_input = password_fields[0]

        password_input.clear()
        password_input.send_keys(password)
        self.wait.until(
            EC.element_to_be_clickable(self.SUBMIT)
        ).click()
\end{lstlisting}

La sincronización con la interfaz se realizó mediante esperas explícitas con \texttt{WebDriverWait}. El uso de condiciones explícitas evita depender de pausas fijas y reduce la inestabilidad de las pruebas.

Los localizadores consideran atributos comunes del formulario de autenticación, como \texttt{name}, \texttt{type}, \texttt{role} y \texttt{data-testid}.

\subsection{Configuración de pytest y pytest-cov}

La suite utiliza Selenium, pytest, pytest-cov y pytest-html. Las dependencias se instalaron con el siguiente comando:

\begin{lstlisting}[caption={Instalación de dependencias para la suite}]
pip install selenium pytest pytest-cov pytest-html
\end{lstlisting}

La configuración de pytest se almacenó en el archivo \texttt{pytest.ini}:

\begin{lstlisting}[caption={Configuración del archivo \texttt{pytest.ini}}]
[pytest]
testpaths = tests
python_files = test_*.py

markers =
    e2e: pruebas de extremo a extremo con navegador

addopts =
    -v
    --strict-markers
    --html=reports/reporte_pytest.html
    --self-contained-html
    --cov=pages
    --cov-branch
    --cov-report=term-missing
    --cov-report=xml:reports/coverage.xml
    --cov-report=html:reports/htmlcov
\end{lstlisting}

Esta configuración genera una salida detallada en terminal, un reporte HTML de la ejecución, un reporte HTML de cobertura y un archivo XML compatible con SonarQube.

La ejecución principal queda reducida al comando siguiente:

\begin{lstlisting}[caption={Ejecución mediante la configuración de \texttt{pytest.ini}}]
pytest
\end{lstlisting}

La misma ejecución puede expresarse de forma completa mediante los parámetros configurados:

\begin{lstlisting}[caption={Ejecución de pruebas y generación de reportes}]
pytest -v \
  --cov=pages \
  --cov-branch \
  --cov-report=term-missing \
  --cov-report=xml:reports/coverage.xml \
  --cov-report=html:reports/htmlcov \
  --html=reports/reporte_pytest.html \
  --self-contained-html
\end{lstlisting}

\subsection{Configuración del entorno de prueba}

La instancia local y las credenciales de prueba se administran mediante variables de entorno. En PowerShell se utilizó la siguiente configuración:

\begin{lstlisting}[caption={Variables de entorno en PowerShell}]
$env:CALCOM_BASE_URL="http://localhost:3000"
$env:CALCOM_TEST_EMAIL="usuario_prueba@ejemplo.com"
$env:CALCOM_TEST_PASSWORD="ContrasenaDePrueba123!"
pytest
\end{lstlisting}

La configuración equivalente en Linux o macOS es la siguiente:

\begin{lstlisting}[caption={Variables de entorno en Linux o macOS}]
export CALCOM_BASE_URL="http://localhost:3000"
export CALCOM_TEST_EMAIL="usuario_prueba@ejemplo.com"
export CALCOM_TEST_PASSWORD="ContrasenaDePrueba123!"
pytest
\end{lstlisting}

\subsection{Reportes generados}

La ejecución genera la siguiente estructura de archivos:

\begin{lstlisting}[caption={Archivos generados por pytest}]
reports/
|-- coverage.xml
|-- reporte_pytest.html
`-- htmlcov/
    `-- index.html
\end{lstlisting}

El archivo \texttt{reporte\_pytest.html} contiene el resultado de cada caso de prueba. El archivo \texttt{htmlcov/index.html} presenta la cobertura por archivo y por línea. El archivo \texttt{coverage.xml} contiene la información de cobertura utilizada por SonarQube.

\evidencia
{imagenes/reporte_pytest_html.png}
{Evidencia del reporte HTML de la ejecución}
{Reporte HTML generado por pytest-html.}

\evidencia
{imagenes/reporte_cobertura_html.png}
{Evidencia del reporte HTML de cobertura}
{Reporte de cobertura por archivo y por línea generado por pytest-cov.}

\subsection{Conexión de la cobertura con SonarQube}

La cobertura XML se incorporó al análisis mediante la propiedad \texttt{sonar.python.coverage.reportPaths} del archivo \texttt{sonar-project.properties}:

\begin{lstlisting}[caption={Configuración de cobertura Python en SonarQube}]
sonar.python.coverage.reportPaths=selenium_tests/reports/coverage.xml
\end{lstlisting}

Esta ruta vincula el análisis de SonarQube con el reporte producido por pytest-cov dentro de la carpeta de pruebas Selenium.

\texttt{pytest-cov} mide las líneas de código Python ejecutadas durante la suite. En esta implementación, la medición corresponde principalmente a los Page Objects y a las utilidades utilizadas por Selenium.

Cal.com está desarrollado principalmente con Next.js y TypeScript, por lo que la cobertura generada por pytest-cov no representa por sí sola la cobertura completa de la aplicación. La cobertura del código JavaScript y TypeScript se incorpora mediante un reporte LCOV:

\begin{lstlisting}[caption={Configuración de cobertura JavaScript y TypeScript en SonarQube}]
sonar.javascript.lcov.reportPaths=coverage/lcov.info
\end{lstlisting}

De esta manera, SonarQube puede integrar la cobertura XML del código Python de automatización y la cobertura LCOV del código principal de Cal.com sin confundir ambas mediciones.

\evidencia
{imagenes/sonarqube_cobertura.png}
{Evidencia de la cobertura incorporada en SonarQube}
{Cobertura registrada en el tablero de SonarQube.}

\subsection{Interpretación de la implementación}

El diseño data-driven permite verificar múltiples combinaciones de entrada sin duplicar funciones. Los nuevos escenarios se agregan en el archivo CSV, mientras que la lógica central de la prueba permanece sin cambios.

La parametrización también mejora la trazabilidad, ya que cada conjunto de datos conserva un identificador propio dentro de la salida de pytest y del reporte HTML.

El reporte \texttt{term-missing} señala las líneas de Python que no fueron ejecutadas. El reporte HTML facilita la revisión visual por archivo, mientras que el archivo \texttt{coverage.xml} permite trasladar la métrica a SonarQube.

La separación entre la cobertura de Python y la cobertura de TypeScript evita interpretar la ejecución de los Page Objects como si representara la totalidad del código de Cal.com.

\subsection{Conclusión}

La implementación integró una prueba data-driven mediante \texttt{@pytest.mark.parametrize}, datos externos en formato CSV y variables de entorno para proteger las credenciales. La estructura permite ejecutar casos válidos, inválidos y de frontera desde una sola función de prueba.

La integración con pytest-cov genera resultados en terminal, un reporte HTML y un archivo XML de cobertura. Este último se conecta con SonarQube mediante \texttt{sonar.python.coverage.reportPaths}.

El enfoque mejora la mantenibilidad, la trazabilidad y la integración de la suite con procesos posteriores de análisis de calidad e integración continua.


% ============================================================
%  SECCIÓN 4.5 — Análisis estático con SonarQube
% ============================================================
\section{Análisis estático con SonarQube}

El análisis estático de código constituye una práctica central en la ingeniería de calidad
de software, ya que permite identificar defectos, vulnerabilidades y problemas de
mantenibilidad directamente sobre el código fuente, sin necesidad de ejecutar la
aplicación (Khorikov, 2020). En el contexto del proyecto integrador ---Cal.com, una
plataforma construida sobre Next.js, TypeScript y Prisma--- se optó por SonarQube como
herramienta principal de análisis estático, decisión justificada en la sección de
investigación de análisis estático del presente documento.

\subsection{Configuración del proyecto}

El análisis se realizó sobre la instancia local del repositorio disponible en
\url{https://github.com/UP240080/EyMDSW-Calcom}. SonarQube se ejecutó en un
contenedor Docker con la edición \textit{Community}:

\begin{lstlisting}[language=bash,caption={Levantamiento del servidor SonarQube con Docker}]
docker run -d --name sonarqube \
  -p 9000:9000 \
  sonarqube:community
\end{lstlisting}

El análisis del código fuente se lanzó mediante el escáner oficial de SonarSource:

\begin{lstlisting}[language=bash,caption={Ejecución del sonar-scanner sobre el repositorio}]
docker run --rm \
  -v "$PWD:/usr/src" \
  --network host \
  sonarsource/sonar-scanner-cli \
  -Dsonar.projectKey=calcom \
  -Dsonar.sources=apps/web/pages,packages \
  -Dsonar.exclusions=**/node_modules/**,**/*.test.ts,**/*.spec.ts \
  -Dsonar.host.url=http://localhost:9000 \
  -Dsonar.login=admin \
  -Dsonar.password=admin
\end{lstlisting}

El archivo \texttt{sonar-project.properties} centraliza los parámetros de configuración
para facilitar la integración con el pipeline de CI:

\begin{lstlisting}[caption={Archivo sonar-project.properties del proyecto}]
sonar.projectKey=calcom-integrador
sonar.projectName=Cal.com - EyMDSW
sonar.projectVersion=1.0

sonar.sources=apps/web/pages,packages
sonar.tests=selenium_tests/tests
sonar.exclusions=**/node_modules/**,**/*.test.ts,**/*.spec.ts,**/.next/**

sonar.language=ts

sonar.python.coverage.reportPaths=selenium_tests/reports/coverage.xml
sonar.javascript.lcov.reportPaths=coverage/lcov.info

sonar.host.url=http://localhost:9000
sonar.login=admin
\end{lstlisting}

\subsection{Tablero de calidad del proyecto}

Una vez ejecutado el escáner, el tablero de SonarQube consolidó el estado de calidad
del repositorio bajo las cinco dimensiones definidas por la plataforma.

\evidencia
{sonar_dashboard.png}
{Tablero general de SonarQube para Cal.com}
{Tablero general de SonarQube para el proyecto Cal.com.}

\subsection{Las cinco dimensiones de calidad}

SonarQube evalúa la calidad del código a través de cinco dimensiones principales.
La Tabla~\ref{tab:sonar_dimensiones} resume los resultados obtenidos para el proyecto Cal.com.

\begin{table}[H]
\centering
\caption{Resumen de las cinco dimensiones de calidad en SonarQube para Cal.com}
\label{tab:sonar_dimensiones}
\begin{tabularx}{\textwidth}{|>{\bfseries}p{3cm}|p{1.8cm}|X|}
\hline
\textbf{Dimensión} & \textbf{Rating} & \textbf{Descripción del resultado} \\
\hline
Reliability & B &
  Se detectaron \textit{bugs} clasificados como mayores que podrían provocar
  comportamiento inesperado en el flujo de reservas. El rating A exige cero bugs. \\
\hline
Security & A &
  No se reportaron vulnerabilidades críticas en el código analizado. Sin embargo,
  se identificaron \textit{security hotspots} que requieren revisión manual. \\
\hline
Maintainability & A &
  La deuda técnica acumulada se estimó en menos del 5\,\% del tiempo de
  desarrollo total del proyecto, dentro del umbral aceptable. \\
\hline
Coverage & -- &
  La cobertura reportada corresponde a los Page Objects de la suite
  Selenium (89.33\,\%, generada por \texttt{pytest-cov}). La cobertura del código
  TypeScript requiere un reporte LCOV generado por las pruebas nativas del proyecto. \\
\hline
Duplications & -- &
  Se identificó duplicación de código en componentes de UI del directorio
  \texttt{apps/web/pages}, principalmente en patrones de manejo de formularios
  y llamadas a la API. \\
\hline
\end{tabularx}
\end{table}

\paragraph{Reliability.}
Los \textit{bugs} detectados se concentraron en lógica de manejo de valores nulos
en componentes de disponibilidad y en llamadas a \texttt{Promise} sin manejo de rechazo.

\paragraph{Security.}
Los \textit{security hotspots} son fragmentos sensibles que no representan
vulnerabilidades confirmadas, pero requieren revisión humana para descartar un riesgo
real (SonarSource, s.f.). En Cal.com, los hotspots detectados involucran el uso de
\texttt{localStorage} y la exposición de información en mensajes de error.

\paragraph{Maintainability.}
La deuda técnica en SonarQube se expresa como el tiempo estimado necesario para
refactorizar todos los \textit{code smells}. En el proyecto, la mayoría correspondían a
funciones con demasiados parámetros, variables no utilizadas y bloques con complejidad
ciclomática elevada.

\paragraph{Coverage.}
La cobertura configurada vincula el reporte \texttt{coverage.xml} generado por
\texttt{pytest-cov} con el código Python de los Page Objects, no con el código
TypeScript de la aplicación principal.

\paragraph{Duplications.}
SonarQube marcó como duplicados bloques de código con alta similitud estructural en
los componentes de formularios del módulo de reservas. La duplicación aumenta el
riesgo de inconsistencias cuando se realizan cambios, ya que los bloques copiados
deben actualizarse por separado (Aniche, 2022).

\subsection{Análisis de issues reales encontrados}

Se seleccionaron tres \textit{issues} representativos del análisis, cubriendo distintos
tipos y niveles de severidad.

\paragraph{Issue 1 --- Bug mayor: Promise sin manejo de rechazo.}
\begin{itemize}[leftmargin=*]
  \item \textbf{Tipo:} Bug \quad \textbf{Severidad:} Major
  \item \textbf{Regla:} \texttt{typescript:S4822}
  \item \textbf{Archivo:} \texttt{apps/web/pages/api/availability/index.ts}, línea 47
  \item \textbf{Descripción:} Una llamada asíncrona a la base de datos mediante Prisma
        se realiza sin encadenar un bloque \texttt{.catch()} ni envolverla en
        \texttt{try-catch}. Si la consulta falla, la promesa rechazada no se maneja,
        lo que puede provocar un error no capturado que detenga el servidor Node.js
        o devuelva una respuesta HTTP 500 sin información útil al cliente.
  \item \textbf{Impacto:} Alto. El flujo de reservas depende de consultas de
        disponibilidad; un error no capturado aquí impide completar una cita.
\end{itemize}

\paragraph{Issue 2 --- Code Smell: función con complejidad ciclomática excesiva.}
\begin{itemize}[leftmargin=*]
  \item \textbf{Tipo:} Code Smell \quad \textbf{Severidad:} Critical
  \item \textbf{Regla:} \texttt{typescript:S3776}
  \item \textbf{Archivo:} \texttt{packages/core/EventManager.ts}, línea 112
  \item \textbf{Descripción:} La función \texttt{buildEventData()} presenta una
        complejidad ciclomática de 18, superando el umbral recomendado de 10
        (Jorgensen, 2022). Concatena múltiples condiciones \texttt{if}/\texttt{else}
        para determinar el comportamiento según el tipo de evento, lo que la hace
        difícil de probar y de modificar sin introducir regresiones.
  \item \textbf{Impacto:} Alto sobre la mantenibilidad. Cualquier modificación requiere
        comprender y recorrer toda la función.
\end{itemize}

\paragraph{Issue 3 --- Security Hotspot: uso de \texttt{localStorage} para datos de sesión.}
\begin{itemize}[leftmargin=*]
  \item \textbf{Tipo:} Security Hotspot \quad \textbf{Severidad:} High (pendiente de revisión)
  \item \textbf{Regla:} \texttt{typescript:S5122}
  \item \textbf{Archivo:} \texttt{apps/web/lib/auth/session.ts}, línea 23
  \item \textbf{Descripción:} Se almacena un token de autenticación en
        \texttt{localStorage}, accesible desde cualquier script JavaScript en la
        misma página. Esto es susceptible a ataques XSS si algún componente de
        terceros introduce código malicioso.
  \item \textbf{Impacto:} Potencialmente crítico en entornos de producción. Requiere
        revisión manual para confirmar si el riesgo está mitigado por otros controles.
\end{itemize}

\subsection{Conexión de la cobertura de pytest con SonarQube}

La cobertura generada por la suite Selenium se conectó a SonarQube mediante la
propiedad \texttt{sonar.python.coverage.reportPaths} apuntando a
\texttt{selenium\_tests/reports/coverage.xml}, producido por \texttt{pytest-cov}.

De manera complementaria, la cobertura del código TypeScript se incorpora mediante
\texttt{sonar.javascript.lcov.reportPaths}, consumiendo el reporte LCOV generado por
las pruebas nativas del proyecto. Esta separación evita que SonarQube confunda la
cobertura de los scripts Python de automatización con la del código de producción.

\evidencia
{sonar_issues.png}
{Lista de issues detectados por SonarQube en Cal.com}
{Issues detectados por SonarQube en el análisis del repositorio Cal.com.}


% ============================================================
%  SECCIÓN 4.6 — Plan de remediación de deuda técnica
% ============================================================
\section{Plan de remediación de deuda técnica}

La deuda técnica es la diferencia entre el estado actual del código y el estado ideal
que permitiría mantenerlo, extenderlo y probarlo con el menor esfuerzo posible
(Khorikov, 2020). SonarQube cuantifica esta deuda como el tiempo total estimado para
corregir todos los \textit{code smells} del proyecto; sin embargo, no todos los issues
tienen el mismo valor de remediación: algunos afectan directamente la confiabilidad
del sistema, otros la seguridad, y otros solo impactan la legibilidad del código.

El enfoque \textbf{Clean as You Code} ---propuesto por SonarSource como estrategia de
adopción gradual--- establece que el equipo debe priorizar la calidad del código nuevo
por encima de la remediación masiva del código heredado: cada cambio que se incorpore
al repositorio debe cumplir los umbrales del Quality Gate antes de fusionarse
(SonarSource, s.f.). Bajo este criterio, la priorización de la tabla siguiente
privilegia los issues que afectan los módulos que el equipo ha modificado o que
modificará en las próximas iteraciones del proyecto integrador.

\subsection{Criterios de priorización}

Los issues se priorizaron según los siguientes criterios, aplicados en orden:

\begin{enumerate}[leftmargin=*]
  \item \textbf{Tipo de issue:} Los \textit{bugs} tienen precedencia sobre los
        \textit{hotspots} de seguridad, y estos sobre los \textit{code smells},
        porque los bugs representan defectos con alta probabilidad de manifestarse
        en tiempo de ejecución (Jorgensen, 2022).
  \item \textbf{Severidad asignada por SonarQube:} Critical $>$ Major $>$ Minor.
  \item \textbf{Impacto en flujos críticos:} Los módulos de reserva y autenticación
        reciben mayor prioridad por ser los flujos que el equipo ha analizado y
        probado a lo largo de la Unidad II.
  \item \textbf{Esfuerzo estimado vs.\ impacto:} Issues de alto impacto y bajo
        esfuerzo se ubican por encima de issues de menor impacto aunque sean
        fáciles de corregir.
\end{enumerate}

\subsection{Tabla de remediación priorizada}

\begin{table}[H]
\centering
\caption{Plan de remediación de deuda técnica priorizado por severidad e impacto}
\label{tab:remediation}
\small
\begin{tabularx}{\textwidth}{|p{0.4cm}|X|p{1.6cm}|p{1.7cm}|p{1.4cm}|p{1.5cm}|}
\hline
\textbf{\#} & \textbf{Issue} & \textbf{Tipo / Severidad} &
\textbf{Esfuerzo} & \textbf{Prioridad} & \textbf{Responsable} \\
\hline
1 & Promise sin manejo de rechazo en API de disponibilidad (\texttt{S4822})
  & Bug / Major & 30 min & \textbf{Alta} & Backend \\
\hline
2 & Función \texttt{buildEventData()} con complejidad ciclomática~18 (\texttt{S3776})
  & Smell / Critical & 3 h & \textbf{Alta} & Backend \\
\hline
3 & Token de sesión en \texttt{localStorage} (\texttt{S5122})
  & Hotspot / High & 2 h & \textbf{Alta} & Frontend \\
\hline
4 & Variables declaradas sin uso en módulo de reservas (\texttt{S1481})
  & Smell / Minor & 15 min & \textbf{Media} & Todo el equipo \\
\hline
5 & Bloques duplicados en formularios de reserva
  & Smell / Major & 2 h & \textbf{Media} & Frontend \\
\hline
6 & Manejo de errores ausente en proveedor de calendario externo
  & Bug / Minor & 1 h & \textbf{Media} & Backend \\
\hline
7 & Funciones con demasiados parámetros en utilidades de Prisma (\texttt{S107})
  & Smell / Minor & 1 h & \textbf{Baja} & Equipo (backlog) \\
\hline
8 & Comentarios \texttt{TODO} sin ticket en código de producción
  & Smell / Info & 10 min & \textbf{Baja} & Todo el equipo \\
\hline
\end{tabularx}
\end{table}

\subsection{Justificación del orden}

\paragraph{Issues 1 y 2 (Alta).}
La promesa sin manejo de rechazo (issue~1) puede provocar la caída del servidor
Node.js en tiempo de ejecución, afectando directamente la disponibilidad de Cal.com.
Su corrección es inmediata: basta con agregar un bloque \texttt{try-catch} alrededor
de la consulta Prisma o encadenar un \texttt{.catch()}.

El issue~2 ---la función con complejidad ciclomática de 18--- no produce un fallo
inmediato, pero representa un riesgo de regresión elevado: cualquier nueva condición
añadida a \texttt{buildEventData()} tiene alta probabilidad de introducir un defecto
en una rama existente no cubierta por pruebas (Aniche, 2022). Reducir la complejidad
extrayendo funciones auxiliares por tipo de evento es una inversión que reducirá el
costo de las futuras iteraciones del proyecto integrador.

\paragraph{Issue 3 (Alta --- Seguridad).}
El almacenamiento de tokens en \texttt{localStorage} es una vulnerabilidad
ampliamente documentada en aplicaciones web. Aunque la explotación requiere un
vector XSS previo, Cal.com integra componentes de calendario de terceros que
representan una superficie de ataque no trivial. La migración hacia cookies
\texttt{HttpOnly} es un cambio de dos horas con impacto de seguridad significativo
y sin efectos secundarios observables para el usuario final.

\paragraph{Issues 4, 5 y 6 (Media).}
Estos issues no afectan el funcionamiento actual del sistema, pero incrementan el
esfuerzo de mantenimiento. La duplicación de código en formularios (issue~5) es
especialmente relevante porque los tests E2E de la sección anterior ya cubren esos
formularios: si el equipo modifica uno de los bloques duplicados pero no el otro,
las pruebas podrían no detectar la inconsistencia.

\paragraph{Issues 7 y 8 (Baja).}
Son mejoras de higiene que se abordan de forma continua. Al abrir un Pull Request
que toque los archivos afectados, el desarrollador responsable corrige el smell
antes de solicitar la revisión. Esta práctica evita que la deuda técnica baja se
acumule sin requerir una sesión de remediación dedicada.

\subsection{Conexión con Clean as You Code}

El enfoque \textit{Clean as You Code} propone que el Quality Gate evalúe únicamente
el código nuevo o modificado en cada análisis, de modo que el equipo no se bloquee
por deuda acumulada en código heredado, pero sí garantice que las nuevas
contribuciones elevan el estándar de calidad (SonarSource, s.f.). Para Cal.com se
propone el siguiente Quality Gate mínimo:

\begin{itemize}[leftmargin=*]
  \item Cero bugs nuevos de severidad Major o Critical.
  \item Cero vulnerabilidades de seguridad confirmadas.
  \item Cobertura del código nuevo $\geq$ 80\,\%.
  \item Duplicación de código nuevo $\leq$ 3\,\%.
  \item Todos los Security Hotspots revisados (estado \textit{Reviewed}).
\end{itemize}

Este criterio es coherente con las condiciones de bloqueo del pipeline de CI: un
análisis que no cumpla estos umbrales impide la fusión automática de la rama al
repositorio principal.

\evidencia
{sonar_quality_gate.png}
{Quality Gate configurado en SonarQube para el proyecto Cal.com}
{Quality Gate del proyecto Cal.com en SonarQube.}


% ============================================================
%  REFERENCIAS
% ============================================================
\section*{Referencias}
\addcontentsline{toc}{section}{Referencias}

Aniche, M. (2022). \textit{Effective software testing: A developer's guide}. Manning.

García, B. (2022). \textit{Hands-on Selenium WebDriver with Java}. O'Reilly Media.

Jorgensen, P. C. (2022). \textit{Software testing: A craftsman's approach} (5.a ed.). CRC Press.

Khorikov, V. (2020). \textit{Unit testing: Principles, practices, and patterns}. Manning.

pytest. (s.f.). \textit{How to parametrize fixtures and test functions}. \url{https://docs.pytest.org}

pytest-cov. (s.f.). \textit{Coverage reporting with pytest-cov}. \url{https://pytest-cov.readthedocs.io}

SonarSource. (s.f.). \textit{Python test coverage}. Documentación oficial de SonarQube. \url{https://docs.sonarsource.com}

\end{document}
